\subsection{Semiconcavity and uniqueness (Evans \S3.3.3)}
\label{sec:3-semiconcavity}
An a.e. PDE alone is insufficient for uniqueness.
For $H(p)=p^2$ and $g=0$ in one dimension, both $u=0$ and
\[
 v(x,t)=\min\{0,|x|-t\}
\]
are Lipschitz and satisfy $v_t+v_x^2=0$ off the three lines
$x=0,\pm t$. At the central upward corner,
$v(h,t)+v(-h,t)-2v(0,t)=2|h|$ for $|h|<t$.
This cannot be bounded by a constant times $h^2$.

\begin{definition}
A function $f$ is semiconcave with constant $C$ if
\[
 f(x+z)-2f(x)+f(x-z)\leq C|z|^2 .
\]
For continuous $f$ this is equivalent to concavity of
$f(x)-C|x|^2/2$, by the midpoint criterion and continuity.
For smooth $f$ it is equivalent to $D^2f\leq CI$.
\end{definition}
\Needspace{45mm}
\PDEAnnotation{半凹只限制向上弯曲}{允许向下尖角，禁止二阶差分过大的向上尖角。它不是要求 Hessian 的绝对值有界，也不是要求函数本身凹；先减去一个二次函数后才凹。}

\begin{theorem}[Source Lemmas 3--4]
If $g$ is semiconcave with constant $C$, the Hopf--Lax function has
the same spatial constant for all $t$.
If instead $D^2H\geq\theta I$, $\theta>0$, it satisfies
\[
 u(x+z,t)-2u(x,t)+u(x-z,t)\leq\frac{|z|^2}{\theta t},
\]
even for initial $g$ that is merely Lipschitz.
\end{theorem}
\begin{expandedproof}[including the dual estimate of Problem 15]
For the first assertion choose a minimizer $y$ at $(x,t)$.
Use $y+z$ and $y-z$ as competitors at $x+z$ and $x-z$.
All three velocities are equal, so their $L$ terms cancel.
The remaining second difference is that of $g$ at $y$.

For the second assertion, uniform convexity makes the maximizer $p(v)$
in $L(v)=\max_p(v\cdot p-H(p))$ unique, with $v=DH(p(v))$.
Strong monotonicity gives
\[
 \theta|p(v)-p(w)|^2
 \leq (v-w)\cdot(p(v)-p(w)),
\]
so $|p(v)-p(w)|\leq|v-w|/\theta$.
Comparing the optimizing expressions at $v$ and $v+h$ shows
$DL(v)=p(v)$; equivalently, apply the implicit function theorem
to $DH(p)=v$. Hence $DL$ is $1/\theta$-Lipschitz.
Integrating its increment along a segment gives
\[
 L(v+h)\leq L(v)+DL(v)\cdot h+\frac{|h|^2}{2\theta}.
\]
Add the estimates for $h$ and $-h$:
$L(v+h)+L(v-h)-2L(v)\leq|h|^2/\theta$.
Now use the same minimizing $y$ at all three target points.
The initial terms cancel and $h=z/t$ gives the claimed constant.

In particular, taking $v=(v_1+v_2)/2$, $h=(v_1-v_2)/2$ proves
\[
 \tfrac12L(v_1)+\tfrac12L(v_2)
 \leq L((v_1+v_2)/2)+\frac{|v_1-v_2|^2}{8\theta}.
\]
This is the dual inequality omitted from the source lemma's proof.
\end{expandedproof}

\begin{theorem}[Uniqueness in the source weak class; source Theorems 7--8]
Let $H$ be smooth, convex and superlinear. Among globally Lipschitz
functions with the same initial datum, satisfying the HJ equation a.e.
and the spatial semiconcavity bound
\[
 u(x+z,t)-2u(x,t)+u(x-z,t)
 \leq C(1+t^{-1})|z|^2,
\]
there is at most one solution.
The Hopf--Lax function belongs to this class if $g$ is semiconcave
or $H$ is uniformly convex.
\end{theorem}
\begin{expandedproof}[the source's regularized cone argument]
Let $w=u-\widetilde u$. Subtract the a.e. PDE and integrate $DH$
on the gradient segment:
\[
 w_t+b\cdot Dw=0,\qquad
 b=\int_0^1 DH(\lambda Du+(1-\lambda)D\widetilde u)\,d\lambda .
\]
Both gradients lie in a fixed ball, so $|b|\leq R$ for a finite $R$.

\paragraph{1. Smooth the coefficients, preserving a one-sided bound.}
Locally convolve $u,\widetilde u$ in space-time and define $b_\varepsilon$
by the same expression with their smoothed spatial gradients.
For $t>2\varepsilon$, semiconcavity implies
\[
 D^2u_\varepsilon,\ D^2\widetilde u_\varepsilon
      \leq C'(1+t^{-1})I .
\]
To check this, integrate their second-difference inequalities against
the nonnegative smoothing kernel, bound $1/s$ by $2/t$ on its time
support, then divide by the squared increment and let it tend to zero.
The gradients remain bounded and converge a.e. to the original ones.

Writing $A=D^2H\geq0$ and $B$ for the convex combination of smoothed
Hessians gives $B\leq C'(1+t^{-1})I$.
Thus $\operatorname{tr}(AB)\leq C'(1+t^{-1})\operatorname{tr}A$,
by conjugating with $A^{1/2}$ and taking traces.
Since $A$ is bounded on the gradient ball,
$\operatorname{div}b_\varepsilon\leq C''(1+t^{-1})$.
Also $|b_\varepsilon|\leq R$ and $b_\varepsilon\to b$ a.e.

\Needspace{45mm}
\PDEAnnotation{凸性把单边 Hessian 界传给速度}{两个半凹函数的 Hessian 只需有上界。因 $D^2H$ 半正定，迹运算仍给出 $\operatorname{div}b$ 的上界；若失去凸性，这一步的符号便没有保证。}
\paragraph{2. Integrate a nonnegative transform in a shrinking cone.}
Take a smooth nonnegative function $\phi$, and put $v=\phi(w)$.
Weak Lipschitz calculus gives $v_t+b\cdot Dv=0$ a.e.
On $B_{\rho(t)}(x_0)$, $\rho(t)=r+R(T-t)$ with $r>0$, set
$e(t)=\int v$. The moving-boundary term and $|b_\varepsilon|\leq R$
make its boundary flux nonpositive. For a.e. $t$ this yields
\[
 e'(t)\leq C''(1+t^{-1})e(t)
       +\int_{B_{\rho(t)}}(b_\varepsilon-b)\cdot Dv.
\]
These integrations hold for Lipschitz $v$ by approximation, or its
Sobolev integration formula on the smooth ball.
For each fixed $\delta>0$, the error tends to zero in $L^1(\delta,T)$
by dominated convergence on the compact cone.
Integrate against a nonnegative time test function and let
$\varepsilon\to0$, obtaining the same differential inequality without
the error on $[\delta,T]$.

\paragraph{3. Start Gronwall away from the singular time zero.}
Let $M$ be the sum of the time Lipschitz constants.
Since $w(\cdot,0)=0$, $|w(x,\delta)|\leq M\delta$.
Choose $\phi=\phi_\delta$ zero on $[-M\delta,M\delta]$ and strictly
positive outside. Then $e(\delta)=0$.
Gronwall on $[\delta,T]$ has a finite coefficient integral, and gives
$e(T)=0$. Hence $|w|\leq M\delta$ on $B_r(x_0)$ at time $T$,
using continuity to pass from a.e. points.
Let $\delta\downarrow0$. The ball and $T$ were arbitrary, so $w=0$.
Existence in the stated class follows from the two semiconcavity
estimates and the preceding Hopf--Lax theorem.
\end{expandedproof}
\Needspace{45mm}
\PDEAnnotation{不能从零时刻硬套 Gronwall}{$1/t$ 在零附近不可积。证明先用时间 Lipschitz 界把差值截掉，在正时刻 $\delta$ 得到零初始能量，再应用 Gronwall，最后令 $\delta\to0$。}

This is the semiconcavity-based weak formulation used in this chapter.
Viscosity solutions provide a broader comparison framework in Chapter 10;
the a.e. equation and the present selection condition must not be
silently interchanged.
